\chapter{Model-Predictive Control and MPC Distillation}
\label{chap:mpc-distill}

\section{Motivation}
\label{sec:mpc-motivation}

\paragraph{What Chapter \ref{chap:offline-online} ruled out.} The diagnostic chain of the previous chapter shows that the LeWM predictor cannot serve as a \emph{source of supervision} for the policy or critic on this benchmark. WM-imagined transitions corrupt TD targets (\S\ref{sec:oo-exp1}); short anchored rollouts do not undo that corruption (\S\ref{sec:oo-exp2}); updating the WM against a critic-derived value loss makes it actively hallucinate high-value states (\S\ref{sec:oo-exp3}); model-uncertainty penalties either fail to discriminate or actively punish task-solving behaviour (\S\ref{sec:oo-exp45}); and analytic gradients through chained flow integration and WM rollouts are too noisy to optimise (\S\ref{sec:oo-exp8}). The encoder is exonerated by the latent-obs control of \S\ref{sec:oo-b2}, so every failure above is structurally a failure of the predictor when used in the training loop.

\paragraph{What it did not rule out.} None of the failure modes touches the much weaker claim that the WM, applied for $1$--$2$ steps starting from a real on-policy latent, can produce an informative \emph{ranking} of candidate action chunks under the learned critic. The WM does not have to be a faithful long-horizon simulator to do that --- it only has to preserve the ordering its rollouts induce over a small candidate set well enough that the highest-scoring action is also a good one in the real environment. The synthesis of \S\ref{sec:oo-synthesis} states this opening explicitly; this chapter exploits it.

\paragraph{The plan.} We continue to train the agent on the offline buffer with the real $(z, \mathbf{a}, r, z')$ transitions, exactly as in Chapter \ref{chap:methodology} --- the critic never sees a WM-imagined next state. What changes is how action chunks are selected. At each decision step we use the WM as a \emph{rollout-and-score oracle}: sample $N$ candidate chunks from the BC-flow actor at the current latent, integrate each one forward in the WM for $H$ steps, and score the resulting trajectory under the critic. The first chunk of the best-scoring candidate is executed. The world model only ever evaluates short horizons starting from a real latent, so the failure mechanisms of Chapter \ref{chap:offline-online} are sidestepped by construction. We then close the loop: whenever this inference-time procedure produces actions that beat the actor's own samples \emph{under the current critic}, we distil those actions back into the actor during continued offline training. The supervision pairs the actor is updated on are $(z, \mathbf{a}^{\text{MPC}})$ where $z$ is real and only the action label is changed; the critic is updated, as before, on the real $(z, \mathbf{a}, r, z')$ tuples from the offline buffer.

\paragraph{Where this sits relative to prior work.} Several established approaches treat the WM as a planner at evaluation time without ever putting it in the TD target:
\begin{itemize}
    \item The LeWM architecture itself \cite{maes2026leworldmodelstableendtoendjointembedding} ships with a CEM-style planner (\S\ref{sec:mpc-bg}) used at evaluation only. Our best-of-$N$ procedure is a learned-proposal variant of that planner, with the BC-flow actor playing the role of the Gaussian sampler in CEM, and we additionally close the loop by distilling MPC-selected actions back into the policy during offline training.
    \item TD-MPC \cite{hansen2022tdmpc} and TD-MPC2 \cite{hansen2024tdmpc2} use a learned model and a learned value function jointly in a sampling-MPC controller, replacing the long planning horizon required by a model alone with a short-horizon look-ahead grounded by the critic. Our rollout-and-score is structurally the same; what we add is the FMQ refinement step (\S\ref{sec:mpc-fmq}), which replaces sampling-and-elite-selection with a closed-form trust-region gradient update.
    \item MOPO \cite{yu2020mopomodelbasedofflinepolicy} and COMBO \cite{yu2021combo} use the WM directly in the TD target with a pessimism penalty. Chapter \ref{chap:offline-online} \S\ref{sec:oo-exp45} already shows that even the strongest version of that idea (NN-distance penalty) collapses to $2.4\%$ on this task, so we deliberately do not put the WM in the TD target.
    \item Dreamer and SVG-style methods \cite{hafner2020dreamer, heess2015svg} use the WM as a differentiable simulator and backpropagate analytic policy gradients through it. \S\ref{sec:oo-exp8} shows this approach collapses on this task.
\end{itemize}
The two methods that work (LeWM-CEM, TD-MPC) and the two that do not (MOPO/COMBO, Dreamer/SVG) sit on opposite sides of the same line: whether the WM appears in the critic's training signal. This chapter stays on the safe side of that line and recovers a $3.2\%$ absolute improvement over the OGBench IMPALA baseline, without ever interacting with the real environment after the offline phase.

\section{Best-of-$N$ MPC at inference}
\label{sec:mpc-bon}

The simplest instantiation of the rollout-and-score procedure is to propose $N$ candidates from the BC-flow actor, evaluate each by rolling it forward $H$ steps in the WM and scoring the resulting trajectory under the critic, and execute the first chunk of the best candidate. We refer to this as \emph{BoN-MPC}. The generic MPC machinery --- proposal sampling, $H$-step rollouts in $f_\psi$, and the choice of score function --- is described in the background (\S\ref{sec:mpc-bg}); here we focus on what is specific to the BoN variant and to our particular instantiation.

\paragraph{Proposal sampling and the connection to Q-chunking.} The proposals are drawn directly from the BC-flow head of the trained FQL actor: sample $N$ noise vectors $\{x_0^{(i)}\}_{i=1}^N \sim \mathcal{N}(0, I_{25})$ and integrate the flow with $S = 10$ Euler steps to obtain $N$ candidate action chunks $\{a_0^{(i)}\}$. This is exactly the best-of-$N$ proposal-selection scheme that Q-chunking introduces for offline-to-online RL \cite{li2025reinforcement}; the only addition we make is to score each candidate not by the critic at the current state but by the critic after an $H$-step WM rollout.

\paragraph{Continuation re-sampling.} When rolling a candidate forward, we draw fresh continuation actions $a_{h+1}, \dots$ from the BC-flow actor at the predicted latents rather than reusing the original candidate's action sequence open-loop. This makes the rollout an unbiased estimate of the on-policy continuation distribution, and it matches how the actor is actually used at deployment: a new chunk is drawn at every WM step, not committed once at the start of planning.

\paragraph{Scoring function.} The full score formulas (Q-term, Q-all, dense$+$Q) are stated in \S\ref{sec:mpc-bg}. The headline numbers in this chapter use \emph{Q-term} scoring --- the critic evaluated only at the terminal $H$-step state --- because empirically this is the most reliable: the critic is trained on real states and is more accurate than the dense reward built on the WM's geometry, especially in the regimes where the WM's calibration is sensitive (\S\ref{sec:oo-exp2}).

\paragraph{Headline result.} Algorithm~\ref{alg:mpc-bon} summarises the BoN-MPC procedure. On the $81.6\%$ offline checkpoint (\texttt{sd55445}), the best configuration ($N = 32$, $H = 1$, Q-term) reaches $88.0\%$ real-env success --- a $+6.4\%$ absolute lift purely from inference-time search, with no parameter updates and no real-env interaction. The full sweep is in Table~\ref{tab:mpc-frozen} in \S\ref{sec:mpc-results}.

\paragraph{Why $H = 1$ wins.} Increasing $H$ to $2$ or $3$ does not help and often hurts: even one extra WM step injects enough prediction error that the critic's evaluation of the terminal state becomes less reliable than its evaluation at the immediate next state. The same trend appears across both scoring functions and is the strongest single argument in this chapter for limiting the WM to its shortest reliable horizon.

\paragraph{The bottleneck.} Standalone BoN-MPC plateaus around $88$--$90\%$ even with larger $N$, because the candidate set is bounded by the BC-flow actor's proposal distribution. To go further we have to either explore actions the actor would not naturally propose --- which is what \S\ref{sec:mpc-grad}--\ref{sec:mpc-fmq} do --- or improve the actor itself so it proposes better candidates --- which is what \S\ref{sec:mpc-interleaved} does.

\begin{algorithm}[t]
\caption{Best-of-$N$ MPC inference}
\label{alg:mpc-bon}
\begin{algorithmic}[1]
\Require current latent $z_t$, BC-flow actor $\pi_\theta$, critic $Q_\theta$, world model $f_\psi$, hyperparameters $N$, $H$, scoring $J$.
\State Tile $z_t$ to a batch $z_{\text{batch}} \in \mathbb{R}^{N \times 192}$.
\State Sample $\{x_0^{(i)}\} \sim \mathcal{N}(0, I_{25})$ and set $a_0^{(i)} \leftarrow \Phi_\theta(z_{\text{batch}}, x_0^{(i)})$.
\State Roll each candidate forward $H$ steps through $f_\psi$ with on-policy continuation re-sampling.
\State Compute score $s_t^{(i)}$ for each candidate using $J$ (Q-term / Q-all / dense$+$Q; see \S\ref{sec:mpc-bg}).
\State \Return $a_0^{(i^\star)}$ with $i^\star = \arg\max_i s_t^{(i)}$.
\end{algorithmic}
\end{algorithm}

\section{Gradient-refined MPC}
\label{sec:mpc-grad}

BoN-MPC is restricted to the actions the proposal distribution happens to sample. Because the WM rollout is differentiable, we can take the best candidate from \S\ref{sec:mpc-bon} and refine it further by gradient ascent on the score function, treating the continuation actions as constants so that the gradient flows only through the initial chunk $a_0$.

\paragraph{Update rule.} Starting from the BoN winner $a_0$, we run $K$ raw gradient ascent steps,
\begin{equation}
    a_0 \;\leftarrow\; \mathrm{clip}\!\Big(a_0 + \eta\,\nabla_{a_0} J(a_0;\,z_t),\;-1,\;1\Big),
    \qquad k = 1, \dots, K,
    \label{eq:mpc-grad}
\end{equation}
with step size $\eta = 0.01$. The continuation noise samples are pre-drawn and fixed across the $K$ steps so that the score being differentiated is well-defined. After the $K$ steps, fresh continuation noises are drawn and a final scoring pass selects the best refined proposal. We refer to this scheme as \emph{Grad-MPC}; it is one of the two refinement variants compared in the FMQ paper \cite{ziakas2026aligning}, and it generalises the standard sampling-MPC controller of TD-MPC \cite{hansen2022tdmpc, hansen2024tdmpc2} by replacing resampling with backpropagation through the WM rollout.

\paragraph{Result.} With $H = 2$, $K = 5$, $N = 32$, Q-term scoring, Grad-MPC reaches $90.0\%$ --- a further $+2.0\%$ over BoN-MPC and now within $2.8\%$ of the IMPALA baseline.

\paragraph{Why Grad-MPC is not the final answer.} Eq.~\eqref{eq:mpc-grad} is a heuristic: $K$ raw gradient steps with a fixed step size, with no built-in trust region. With small $K$ the refinement is too gentle; with large $K$ it overshoots into regions where the critic has not been trained. The right setting of $\eta$ and $K$ is task-dependent and finicky --- a sweep over $K \in \{5, 8, 10, 15\}$ in Table~\ref{tab:mpc-fmq-distilled} shows non-monotonic behaviour, with $K = 8$ matching the best FMQ result at one operating point and $K = 10$ already degrading. The next section replaces this heuristic with a closed-form update that is, in a precise sense, the unique optimum.

\section{Flow Map $Q$-Guidance (FMQ)}
\label{sec:mpc-fmq}

The conceptual move is to formulate the action-refinement problem as a trust-region optimisation problem and ask for the velocity update that maximises the linear approximation of the critic under a Euclidean constraint. The closed-form solution is what \cite{ziakas2026aligning} call \emph{Flow Map Q-Guidance} (FMQ), and we use it as a one-step optimal alternative to the $K$-step heuristic of Grad-MPC.

\paragraph{Trust-region objective.} Let $\pi^{\mathrm{ref}}(\cdot \mid z)$ be a reference flow-map policy producing actions $a_1 = a_r + (1-r)\,u^{\mathrm{ref}}_{r,1}(a_r \mid z)$ through the two-time average velocity $u^{\mathrm{ref}}_{r,1}$ (cf.\ \S\ref{sec:fmq-bg}). We seek the velocity perturbation that maximises the first-order expansion of $Q_\phi$ around $a_1$ subject to a Euclidean ball of radius $\eta$ in velocity space:
\begin{equation}
    \max_{u_{r,1}}\; Q_\phi\!\big(z,\, a_r + (1-r)\, u_{r,1}(a_r \mid z)\big)
    \qquad\text{s.t.}\qquad
    \big\|u_{r,1} - u^{\mathrm{ref}}_{r,1}\big\|_2 \;\le\; \eta.
    \label{eq:mpc-fmq-objective}
\end{equation}
Theorem 3.2 of \cite{ziakas2026aligning} solves this in closed form. Substituting the flow-map parametrisation into the first-order Taylor expansion of $Q_\phi$ and applying the KKT conditions to the constrained linearised problem yields a unique optimal perturbation: the normalised $Q$-gradient at $a_1$.
\begin{equation}
    \boxed{\;u^{\,*}_{r,1}(a_r \mid z) \;=\; u^{\mathrm{ref}}_{r,1}(a_r \mid z) \;+\; \eta\,\frac{\nabla_a Q_\phi(z, a_1)}{\big\|\nabla_a Q_\phi(z, a_1)\big\|_2}\;.}
    \label{eq:mpc-fmq-target}
\end{equation}
Two properties of this expression are worth emphasising. First, it is \emph{the} closed-form optimum of Eq.~\eqref{eq:mpc-fmq-objective} rather than one of several reasonable choices: any other perturbation within the $\eta$-ball can only do worse in the local linear approximation. Second, the magnitude of the update is fixed by the radius $\eta$ rather than by the magnitude of the gradient; the direction is the normalised gradient. This makes FMQ insensitive to the scale of $\nabla Q_\phi$ in a way that Grad-MPC's $K$ raw steps are not.

\paragraph{Inference-time FMQ.} Used as a planner, FMQ proceeds exactly like Grad-MPC except that the inner ascent loop of Eq.~\eqref{eq:mpc-grad} is replaced by a single application of Eq.~\eqref{eq:mpc-fmq-target}. The two methods share the BoN proposal sampling and the score function; they differ only in the refinement step. Algorithm~\ref{alg:mpc-refine} presents both as branches of the same procedure to make the relationship explicit.

\begin{algorithm}[t]
\caption{Gradient-refined MPC inference. Grad-MPC ($K$ raw steps) and FMQ (one normalised step) share the proposal-and-score machinery and differ only in the inner refinement.}
\label{alg:mpc-refine}
\begin{algorithmic}[1]
\Require $z_t$, $\pi_\theta$, $Q_\theta$, $f_\psi$, scoring $J$; refinement-specific hyperparameters: $(K, \eta_{\text{Grad}})$ for Grad-MPC, $\eta$ for FMQ.
\State Sample initial proposals $\{a_0^{(i)}\}_{i=1}^N \sim \pi_\theta(\cdot \mid z_t)$ via $N$ Euler integrations of the BC flow.
\State Pre-sample continuation noises $\{\tilde x_0^{(i, h)}\}$, $h = 0, \dots, H-1$, held constant during refinement.
\If{Grad-MPC}
    \For{$k = 1, \dots, K$}
        \State $g \leftarrow \nabla_{a_0}\, J(a_0;\,z_t,\,\{\tilde x_0\})$ via BPTT through the $H$-step WM rollout.
        \State $a_0 \leftarrow \mathrm{clip}\big(a_0 + \eta_{\text{Grad}}\, g,\;-1,\;1\big)$.
    \EndFor
\ElsIf{FMQ}
    \State $g \leftarrow \nabla_{a_0}\, J(a_0;\,z_t,\,\{\tilde x_0\})$ via a single BPTT pass.
    \State $a_0 \leftarrow \mathrm{clip}\big(a_0 + \eta\, g / \|g\|_2,\;-1,\;1\big)$.
\EndIf
\State Draw fresh continuation noises and re-score $a_0^{(1)}, \dots, a_0^{(N)}$ as in Algorithm~\ref{alg:mpc-bon}.
\State \Return $a_0^{(i^\star)}$ with $i^\star = \arg\max_i s_t^{(i)}$.
\end{algorithmic}
\end{algorithm}

\paragraph{Sensitivity on the base critic.} The FMQ-optimal direction in Eq.~\eqref{eq:mpc-fmq-target} only matters if the critic's value estimate in that direction is reliable. The critic of the offline FQL agent was trained on real $(z, a, r, z')$ tuples in which $a$ was sampled from the BC-flow actor. FMQ-shaped actions --- actions displaced along the normalised $Q$-gradient --- are not in the distribution the critic has seen. Empirically (Table~\ref{tab:mpc-fmq-basepolicy}), the base critic peaks under FMQ at $\eta \approx 0.1$ and collapses sharply for $\eta > 0.1$, reaching only $57.2\%$ at $\eta = 0.3$. Large normalised steps move the action into a region of unreliable $Q$, the score ordering collapses, and FMQ is no better than random sampling. This is the empirical motivation for \S\ref{sec:mpc-interleaved}: if we train the critic on FMQ-shaped actions, it learns to be reliable on the action distribution FMQ produces, and we can use much finer trust regions at inference.

\paragraph{FMQ as a training-time learning target.} Eq.~\eqref{eq:mpc-fmq-target} also serves a second role beyond inference. If we set the reference velocity $u^{\mathrm{ref}}_{r,1}$ equal to the frozen offline flow-map velocity $u^{\text{off}}_{r,1}$, the right-hand side becomes a closed-form self-bootstrapped learning target for the online flow-map velocity $u^\theta_{r,1}$: a velocity that is exactly $\eta$ further along the $Q$-gradient than the offline reference. We regress against this target with a stop-gradient on the right-hand side,
\begin{equation}
    \mathcal{L}_{\text{FMQ}}(\theta) \;=\; \mathbb{E}_{r, a_0, \mathbf{a}_{\text{data}}}\!\left[\,
        \Big\| u^\theta_{r,1}(a_r \mid z) \;-\; \mathrm{sg}\!\Big(u^{\text{off}}_{r,1}(a_r \mid z) \;+\; \eta\,\hat g_\phi\Big) \Big\|_2^2\,\right],
    \quad \hat g_\phi = \frac{\nabla_a Q_\phi(z, a_1)}{\|\nabla_a Q_\phi(z, a_1)\|_2 + \kappa},
    \label{eq:mpc-fmq-loss}
\end{equation}
with a small $\kappa$ for numerical stability, and using a $u$-time interpolant $a_r = (1 - r)\,a_0 + r\,\mathbf{a}_{\text{data}}$ between noise $a_0 \sim \mathcal{N}(0, I)$ and an offline action $\mathbf{a}_{\text{data}}$. This is Eq.~(12) of the FMQ paper. We use it as the actor-side label generator inside the interleaved distillation loop of \S\ref{sec:mpc-interleaved}, and it is the source of the $96\%$ headline result. The FMQ paper additionally introduces an adaptive per-sample $\eta$ derived from critic-ensemble disagreement; we use the same fixed value $\eta = 0.3$ as their baseline.

\section{Offline MPC distillation}
\label{sec:mpc-offline-distill}

The simplest way to fold the inference-time gains of \S\S\ref{sec:mpc-bon}--\ref{sec:mpc-fmq} into the policy itself is a two-phase pipeline: relabel the offline dataset with MPC-selected actions in a one-shot pass, then continue training the BC-flow actor on those relabels.

\paragraph{Phase A --- relabel.} For every latent observation $z \in \mathcal{Z}_{\text{off}}$ in the offline cache we run the Grad-MPC procedure of \S\ref{sec:mpc-grad} ($N = 32$, $H = 2$, $K = 5$, $\eta = 0.01$, Q-term scoring) using the offline-trained actor and critic and the fully fine-tuned WM, and record the resulting chunk as $\mathbf{a}^{\text{MPC}}(z)$. The output is a relabelled dataset of $(z, \mathbf{a}^{\text{MPC}}(z))$ pairs.

\paragraph{Phase B --- distil.} We continue training only the BC-flow actor on the relabelled pairs, with the critic and one-step distillation heads frozen. The loss is the flow-matching BC loss of Eq.~\eqref{eq:bc-flow-bg} with $\mathbf{a}$ replaced by $\mathbf{a}^{\text{MPC}}$:
\begin{equation}
    \mathcal{L}_{\text{distil}}(\theta) \;=\;
        \mathbb{E}_{z, t, x_0}\!\left[\,
            \big\| u_\theta(z, x_t, t) - (\mathbf{a}^{\text{MPC}}(z) - x_0) \big\|_2^2\,\right],
    \label{eq:mpc-distil-loss}
\end{equation}
with $x_t = (1-t)\,x_0 + t\,\mathbf{a}^{\text{MPC}}(z)$. The optimiser is AdamW with learning rate $10^{-4}$, weight decay $10^{-4}$, batch size $256$, for $10^5$ gradient steps.

\paragraph{Phase C --- evaluate.} The distilled actor can be evaluated standalone (BoN $N = 4$ under the original critic, no WM at test time) or used as the proposal distribution inside Grad-MPC or FMQ.

\paragraph{The limitation of one-shot distillation.} Standalone performance does lift modestly --- to $\sim 81$--$84\%$ depending on hyperparameters --- but no further. The deeper issue is structural: the labels in Phase A are generated by a fixed actor and a fixed critic, and once they are baked into the dataset they are stale for the rest of training. As the actor improves under $\mathcal{L}_{\text{distil}}$, the labels are no longer optimal under the improved proposal distribution; as the critic could be improved if it were also being updated, the scoring function used to generate the labels is also stale. The interleaved variant in the next section fixes both by refreshing the labels periodically using the current actor and critic.

\section{Interleaved MPC distillation}
\label{sec:mpc-interleaved}

The MPC labels of \S\ref{sec:mpc-offline-distill} depend on both the current actor (which proposes candidates) and the current critic (which scores them). If we let the actor improve under those labels and the critic improve alongside it on real offline transitions, the labels themselves become better on the next refresh --- the actor proposes better candidates and the critic scores them more reliably. This is a positive-feedback loop. The one-shot pass of \S\ref{sec:mpc-offline-distill} discards it; the interleaved variant of this section preserves it.

\paragraph{Mechanism.} Inside the standard offline FQL training loop on $\mathcal{D}_{\text{off}}$ (\S\ref{sec:oo-setup}), we maintain a circular buffer $\mathcal{B}_{\text{MPC}}$ of $(z, \mathbf{a}^{\text{MPC}}(z))$ pairs. After a warm-up of $T_{\text{warm}}$ steps, the buffer is refreshed every $T_{\text{relabel}}$ gradient steps: $|\mathcal{B}_{\text{MPC}}|$ states are sampled uniformly from the offline buffer and the MPC procedure --- either Grad-MPC or FMQ --- is run on each using the \emph{current} actor and critic. At every training step we draw a real minibatch from the offline buffer and replace the first $\lfloor pB \rfloor$ rows with samples from $\mathcal{B}_{\text{MPC}}$, then run the standard FQL update on the resulting mixed batch.

\paragraph{Critically, the critic still sees only real transitions.} The replacement happens to both the observations and the actions in the affected rows, but the rewards and next-states stay anchored to the real offline tuples. The actor is supervised on $(z, \mathbf{a}^{\text{MPC}})$ where $z$ is a real latent; the critic is updated with the real reward and next-state stored alongside that latent in the offline dataset. The WM never appears in the critic's TD target.

\paragraph{Mixed vs separate injection.} An alternative scheme runs the MPC-label update as an actor-only step on top of the main FQL update, leaving the critic's batch composition untouched. Our experiments show that the mixed mode wins decisively (\S\ref{sec:mpc-results}). The reason is that the critic must \emph{see} MPC actions to assign them high $Q$ --- if the critic is updated only on actor samples while the actor is being pulled toward MPC labels, the actor is dragged in a direction the critic cannot endorse, and the loop fails to close.

\paragraph{Grad-MPC labels vs FMQ labels.} The label generator inside Algorithm~\ref{alg:mpc-interleaved} can be either Grad-MPC (\S\ref{sec:mpc-grad}) or FMQ (\S\ref{sec:mpc-fmq}); the rest of the training loop is unchanged. Grad-MPC labels produce a strong but not optimal result: the headline Grad-MPC-distilled run reaches $92.4\%$ on MPC-eval. FMQ labels --- being the closed-form optimum of the trust-region problem of Eq.~\eqref{eq:mpc-fmq-objective} --- produce qualitatively cleaner action displacements that lie closer to actions the actor can match, and the critic adapts to them more cleanly. The FMQ-distilled run with $T_{\text{relabel}} = 50\,000$ reaches a standalone peak of $86.8\%$ at $400$k steps and, evaluated with FMQ inference at the fine inference radius $\eta = 0.02$, reaches \textbf{$96.0\%$} (Table~\ref{tab:mpc-fmq-distilled}), $+3.2\%$ over the IMPALA baseline.

\paragraph{Optional advantage filter.} An AWR-style filter that retains only labels for which $Q_\phi(z, \mathbf{a}^{\text{MPC}}) > Q_\phi(z, \pi_\theta(z))$ is straightforward to add and is included as an ablation. Empirically it gives no consistent improvement over plain mixed injection on this task (Table~\ref{tab:mpc-interleaved}), so we report it for completeness but use the plain variant as the default.

\begin{algorithm}[t]
\caption{Interleaved MPC distillation.}
\label{alg:mpc-interleaved}
\begin{algorithmic}[1]
\Require offline dataset $\mathcal{D}_{\text{off}}$, FQL agent $(\pi_\theta, Q_\phi)$, frozen WM $f_\psi$, hyperparameters $T_{\text{warm}}$, $T_{\text{relabel}}$, buffer size $|\mathcal{B}_{\text{MPC}}|$, mix ratio $p$, label generator $\mathcal{M} \in \{\text{Grad-MPC},\, \text{FMQ}\}$.
\State $\mathcal{B}_{\text{MPC}} \leftarrow \emptyset$.
\For{$\text{step} = 1, \dots, T_{\text{total}}$}
    \If{$\text{step} \ge T_{\text{warm}}$ \textbf{and} $\text{step} \bmod T_{\text{relabel}} = 0$}
        \State Sample $|\mathcal{B}_{\text{MPC}}|$ states from $\mathcal{D}_{\text{off}}$.
        \State For each $z$, set $\mathbf{a}^{\text{MPC}}(z) \leftarrow \mathcal{M}(z;\,\pi_\theta,\,Q_\phi,\,f_\psi)$ (Alg.~\ref{alg:mpc-refine}).
        \State $\mathcal{B}_{\text{MPC}} \leftarrow \{(z,\,\mathbf{a}^{\text{MPC}}(z))\}$; reset the actor-only optimiser state.
    \EndIf
    \State Sample minibatch $B$ from $\mathcal{D}_{\text{off}}$.
    \If{$\mathcal{B}_{\text{MPC}} \ne \emptyset$}
        \State Sample $\lfloor pB \rfloor$ pairs from $\mathcal{B}_{\text{MPC}}$ and overwrite the first $\lfloor pB \rfloor$ rows of $B$'s observations and actions.
    \EndIf
    \State Update $(\pi_\theta, Q_\phi)$ on $B$ with the standard FQL critic and actor losses (Eqs.~\eqref{eq:bc-flow-bg}--\eqref{eq:fql-critic}).
\EndFor
\end{algorithmic}
\end{algorithm}

\section{Results}
\label{sec:mpc-results}

We collect the headline numbers here; the full hyperparameter sweeps are in Appendix~\ref{app:oo-sweeps}.

\begin{table}[h]
\centering
\caption{Frozen offline policy combined with WM-based inference
   (Algs.~\ref{alg:mpc-bon}--\ref{alg:mpc-refine}). All entries use
   the $81.6\%$ offline checkpoint (\texttt{sd55445}); $N = 32$ throughout, $250$ episodes on
   \texttt{visual-cube-single-singletask-task1-v0}.}
\label{tab:mpc-frozen}
\begin{tabular}{l c c l r}
\toprule
Method & $H$ & $K$ & Scoring & SR (\%) \\
\midrule
BoN-MPC                 & 1 & 0 & dense$+$Q & 87.2 \\
BoN-MPC                 & 1 & 0 & Q-term    & \textbf{88.0} \\
BoN-MPC                 & 2 & 0 & Q-term    & 83.6 \\
BoN-MPC                 & 3 & 0 & dense$+$Q & 81.6 \\
\midrule
Grad-MPC                & 1 & 5 & Q-term    & 88.8 \\
Grad-MPC                & 2 & 5 & Q-term    & \textbf{90.0} \\
Grad-MPC                & 3 & 5 & dense$+$Q & 88.8 \\
\midrule
BoN-MPC (VAML-tuned WM) & 3 & 0 & Q-term    & 81.6 \\
\bottomrule
\end{tabular}
\end{table}

\begin{table}[h]
\centering
\caption{FMQ inference (Eq.~\eqref{eq:mpc-fmq-target}) applied to the
   $81.6\%$ frozen offline checkpoint. The base critic peaks under FMQ
   at $\eta = 0.1$ and collapses sharply above that; this sensitivity
   motivates the interleaved FMQ-label training in \S\ref{sec:mpc-interleaved}.
   $N = 32$, $H$ as shown.}
\label{tab:mpc-fmq-basepolicy}
\begin{tabular}{c l c r}
\toprule
$H$ & Scoring & $\eta$ & SR (\%) \\
\midrule
2 & Q-all  & 0.10 & \textbf{90.0} \\
2 & Q-all  & 0.30 & 57.2 \\
2 & Q-all  & 0.50 & 25.6 \\
1 & Q-all  & 0.30 & 60.0 \\
2 & Q-term & 0.30 & 58.8 \\
1 & Q-term & 0.30 & 61.2 \\
\bottomrule
\end{tabular}
\end{table}

\begin{table}[h]
\centering
\caption{Interleaved MPC distillation with Grad-MPC labels (Alg.~\ref{alg:mpc-interleaved}, $\mathcal{M} = \text{Grad-MPC}$, $K = 5$).
   $500\,000$ offline steps, seed $42$, starting from the $81.6\%$ offline checkpoint. \emph{Standalone} is the best per-checkpoint success rate
   when the actor is used without the WM (BoN $N = 4$); \emph{MPC-eval} re-applies Grad-MPC ($N = 32$, $H = 2$, $K = 5$) on top of the final actor.
   Headline configuration (full, mixed, $p = 0.3$) reaches $\mathbf{92.4\%}$ MPC-eval.}
\label{tab:mpc-interleaved}
\begin{tabular}{l l l c r r}
\toprule
Configuration & MPC HP & Mode & $p$ & Peak SA (\%) & MPC-eval (\%) \\
\midrule
Grad-MPC, light, actor-only  & $N{=}8$, $H{=}1$, $K{=}3$   & sep   & 1.0 & 83.6 & 86.8 \\
Grad-MPC, light, mixed       & $N{=}8$, $H{=}1$, $K{=}3$   & mixed & 0.3 & 81.6 & 91.2 \\
Grad-MPC, full, actor-only   & $N{=}32$, $H{=}2$, $K{=}5$  & sep   & 1.0 & 84.0 & 85.6 \\
Grad-MPC, full, mixed        & $N{=}32$, $H{=}2$, $K{=}5$  & mixed & 0.3 & 82.8 & \textbf{92.4} \\
AWR-filtered, actor-only     & $N{=}32$, $H{=}2$, $K{=}5$  & awr   & 0.3 & 84.4 & 90.0 \\
AWR-filtered, mixed          & $N{=}32$, $H{=}2$, $K{=}5$  & mixed-awr & 0.3 & 84.4 & 91.2 \\
\bottomrule
\end{tabular}
\end{table}

\begin{table}[h]
\centering
\caption{FMQ-label interleaved distillation: relabelling interval sweep. All runs use FMQ labels at $\eta_{\text{train}} = 0.1$,
   $H = 2$, $N = 32$, mixed injection at $p = 0.3$. \emph{Grad-MPC labels at 25 k} reuses the Grad-MPC label recipe at a $25$k interval for direct comparison. MPC-eval uses Grad-MPC ($N = 32$, $H = 2$, $K = 5$, Q-term).}
\label{tab:mpc-relabel-sweep}
\begin{tabular}{l l l r r r r r}
\toprule
& & & \multicolumn{5}{c}{SR at checkpoint (\%)} \\
\cmidrule(lr){4-8}
Run & Labels & Interval & 100k & 200k & 300k & 400k & 500k \\
\midrule
\multicolumn{8}{l}{\emph{Standalone (BoN $N = 4$)}} \\
FMQ-distil, 25k & FMQ      & 25k & 67.2 & \textbf{86.8} & 80.0 & 82.0 & 71.2 \\
FMQ-distil, 50k & FMQ      & 50k & 68.4 & 80.8           & 68.0 & \textbf{86.8} & 74.0 \\
\midrule
\multicolumn{8}{l}{\emph{MPC-eval (Grad-MPC $K = 5$, Q-term)}} \\
FMQ-distil, 25k & FMQ      & 25k & 80.8 & 88.4 & 91.2 & 91.6 & 88.8 \\
FMQ-distil, 50k & FMQ      & 50k & 79.2 & 88.8 & 79.2 & \textbf{93.2} & 90.4 \\
Grad-MPC labels at 25k & Grad-MPC & 25k & 62.8 & 74.4 & 67.6 & 79.2 & --   \\
\bottomrule
\end{tabular}
\end{table}

\begin{table}[h]
\centering
\caption{Inference-time sweep on the FMQ-distilled actor (the 400k checkpoint of \emph{FMQ-distil, 50k}). The distilled critic is robust to small FMQ radii: SR stays above $94\%$
   for $\eta \in [0.01, 0.05]$ and reaches $96.0\%$ at $\eta \in \{0.02, 0.03\}$, a $+3.2\%$ lift over the IMPALA baseline. Compare with
   Table~\ref{tab:mpc-fmq-basepolicy}, where the un-distilled critic collapses at any $\eta > 0.1$.}
\label{tab:mpc-fmq-distilled}
\begin{tabular}{l l l r}
\toprule
Method & Scoring & Config & SR (\%) \\
\midrule
FMQ      & Q-all  & $\eta = 0.01$ & 94.8 \\
FMQ      & Q-all  & $\eta = 0.02$ & \textbf{96.0} \\
FMQ      & Q-all  & $\eta = 0.03$ & \textbf{96.0} \\
FMQ      & Q-all  & $\eta = 0.05$ & 94.0 \\
FMQ      & Q-all  & $\eta = 0.10$ & 90.8 \\
FMQ      & Q-all  & $\eta = 0.20$ & 80.8 \\
FMQ      & Q-term & $\eta = 0.05$ & 93.6 \\
\midrule
Grad-MPC & Q-term & $K = 5$       & 93.2 \\
Grad-MPC & Q-term & $K = 8$       & \textbf{96.0} \\
Grad-MPC & Q-term & $K = 10$      & 90.8 \\
\midrule
Standalone (BoN $N = 4$) & -- & -- & 86.8 \\
IMPALA baseline          & -- & -- & 92.8 \\
\bottomrule
\end{tabular}
\end{table}

\paragraph{Headline progression.}
\begin{enumerate}
    \item Offline baseline (\texttt{sd55445}): $81.6\%$.
    \item Frozen actor + BoN-MPC ($H = 1$, $N = 32$): $88.0\%$.
    \item Frozen actor + Grad-MPC ($H = 2$, $K = 5$): $90.0\%$.
    \item Interleaved Grad-MPC distillation (full, mixed): $92.4\%$, first method within $0.4\%$ of the IMPALA reference.
    \item Interleaved FMQ-label distillation + Grad-MPC inference (FMQ-distil 50k, 400k checkpoint, $K = 5$): $93.2\%$, first to beat IMPALA.
    \item \textbf{Interleaved FMQ-label distillation + FMQ inference} (same checkpoint, $\eta = 0.02$): \textbf{$96.0\%$}.
\end{enumerate}

\paragraph{No real-env interaction in the post-offline phase.} The $96.0\%$ result beats the IMPALA reference --- itself trained from privileged state with millions of real-env steps --- using only the offline buffer and the frozen LeWM. The agent never collects a real-env transition after the initial offline pretraining; the entire $+14.4\%$ absolute improvement over the offline baseline is extracted from inference-time WM rollouts and the closed loop that distils those rollouts back into the actor.

Further hyperparameter sweeps (FMQ label radius, relabelling interval, mix ratio, buffer size, actor learning rate, training duration) and the $\eta = 0.02$ evaluation on the best-of-each-run checkpoints are deferred to Appendix~\ref{app:oo-sweeps}.

\section{Offline-to-online inside the world model}
\label{sec:mpc-online-wm}

Every method up to this point has been careful to keep the world model out of the critic's training signal. The two distillation schemes of \S\S\ref{sec:mpc-offline-distill}--\ref{sec:mpc-interleaved} change only the \emph{action} the actor is supervised on; the critic continues to bootstrap on the real offline tuples $(z, \mathbf{a}, r, z')$. This section asks the question those schemes deliberately avoid: \emph{can we reintroduce world-model transitions into the buffer at all?} This is exactly the offline-to-online-inside-the-WM setup that \S\ref{sec:oo-exp1} showed collapses --- the agent acts in the WM, gets imagined feedback, and trains on it. We show that with one structural change it not only stops collapsing but recovers the strongest standalone policy of the whole chapter, and we explain precisely why.

\paragraph{The hidden inconsistency in action-only distillation.} The interleaved scheme of \S\ref{sec:mpc-interleaved} overwrites the action of a real transition with the MPC-selected chunk and leaves the stored reward and next-state untouched, so the critic is updated on $(z,\, \mathbf{a}^{\text{MPC}},\, r,\, z')$. We argued there that this keeps the WM out of the TD target, which is true; but it introduces a different defect. The next-state $z'$ stored alongside $z$ is the consequence of the \emph{data} action $\mathbf{a}$, not of $\mathbf{a}^{\text{MPC}}$. Taking $\mathbf{a}^{\text{MPC}}$ from $z$ would, in general, land somewhere else entirely. The tuple the critic is asked to fit therefore describes a transition that never happened: it pairs an action with the reward and successor of a \emph{different} action. The loop closes (the critic does learn to assign high $Q$ to $\mathbf{a}^{\text{MPC}}$, which is why the mixed scheme outperforms the actor-only one) but it closes on a dynamically inconsistent target, and the only thing preventing this from corrupting the value function is that the displacement $\mathbf{a}^{\text{MPC}} - \mathbf{a}$ is small enough that $z'$ is still approximately right.

\paragraph{World-model-consistent transitions.} The principled repair is to stop reusing the stale successor and instead ask the WM what $\mathbf{a}^{\text{MPC}}$ actually leads to. We run the agent as a closed-loop controller \emph{inside the world model}: from a real initial latent we select a chunk by FMQ-MPC (\S\ref{sec:mpc-fmq}), step the WM to obtain the imagined successor $z'_\psi = f_\psi(z, \mathbf{a}^{\text{MPC}})$, assign a sparse goal-reaching reward $r_\psi = \mathbb{1}\{\lVert z'_\psi - z_{\text{goal}}\rVert < \tau\}$, and continue the rollout from $z'_\psi$ until success or a $40$-step truncation. Each step contributes a transition $(z,\, \mathbf{a}^{\text{MPC}},\, r_\psi,\, z'_\psi)$ that is now \emph{self-consistent} --- the successor really is the model's prediction for the action taken --- and these are added to a replay buffer seeded with the offline data. The cost of consistency is that $z'_\psi$ is imagined: we have traded the action-mismatch of \S\ref{sec:mpc-interleaved} for a direct dependence on the quality of the WM. This is the dependence Chapter~\ref{chap:offline-online} spent its length warning against.

\paragraph{Two trusts that make the dependence safe.} The danger is real, but it is bounded by two assumptions, each of which is much weaker than ``the WM is a faithful simulator,'' and each of which the diagnostics of the previous chapters actually support.
\begin{enumerate}
    \item \textbf{We trust states our value function can rank.} The chunk that is rolled out is not an arbitrary action --- it is the $\arg\max$ over a candidate set under the \emph{offline} critic, which was trained only on real $(z,\mathbf{a},r,z')$ and is therefore reliable on the data manifold. The critic acts as an in-distribution gate: a candidate whose imagined trajectory wanders off the manifold is scored by a $Q$ the critic cannot vouch for and is rejected in favour of one that stays on it. The WM is never required to be globally accurate, only locally rank-preserving over a handful of one- to two-step rollouts --- precisely the weak claim left open by the synthesis of \S\ref{sec:oo-synthesis}.
    \item \textbf{We trust the WM to rank a small set of actions.} Using the model to choose between a few candidate action chunks over a short horizon is exactly how the LeWM architecture uses it in the original paper --- its CEM planner (\S\ref{sec:mpc-bg}) at evaluation time \cite{maes2026leworldmodelstableendtoendjointembedding}. We are not asking the WM to do anything it was not built to do; we are reusing its sanctioned mode of operation as the data-generation mechanism for an online phase.
\end{enumerate}
The combination is what distinguishes this setup from the failure of \S\ref{sec:oo-exp1}: there, raw WM rollouts entered the buffer with no value-based gate on which states were admissible, and the critic bootstrapped on all of them; here, the clean offline critic both \emph{selects} the actions that generate the data and is kept frozen so it never has to bootstrap on the imagined successors.

\paragraph{Experimental setup (E2).} We test this on a single task, \texttt{visual-cube-single-singletask-task1-v0}, using the fully fine-tuned LeWM. The starting point is a fresh pure-offline FQL checkpoint (seed $42$) whose standalone success is $76.4\%$ at the final step and $88.0\%$ at its peak. The online phase runs $500\,000$ gradient steps; at each step one FMQ-MPC action (full config: $N = 32$, $H = 2$, $\eta = 0.02$, Q-all scoring with one FMQ refinement step) is selected and rolled in the WM, one imagined transition is appended to the buffer, and a minibatch is drawn from the combined offline-plus-imagined buffer. Crucially the WM and the planner are used \emph{only} to generate training data: evaluation is standalone (best-of-$4$ under the critic, no planner, no WM at test time), $250$ episodes, reported as mean $\pm$ std over three seeds. The agent never touches the real environment during the online phase.

\paragraph{The critic must not bootstrap on imagined transitions (E2a vs.\ E2c).} The two runs in Table~\ref{tab:mpc-online-e2} differ in exactly one bit. In \textbf{E2a} the standard FQL update is applied to the mixed buffer, so the critic bootstraps on the imagined transitions; the result is a collapse to $30.7 \pm 20.8\%$, with individual seeds finishing at $60$, $14$, and $18\%$ after transient early peaks near $90\%$. Two mechanisms compound here. The first is the value-corruption mechanism of \S\ref{sec:oo-exp3} --- the WM in the TD target. The second is a reward-convention clash: the offline OGBench transitions carry the native cost-to-goal reward ($r = -1$ per step until success, $Q \in [-100, 0]$) \cite{ogbench_park2025}, whereas the imagined transitions carry the goal-reaching reward $r_\psi \in \{0, 1\}$ ($Q \in [0, 1]$); a single critic cannot fit two incompatible value scales in one buffer. In \textbf{E2c} we freeze the critic for the entire online phase --- only the BC-flow actor is trained, distilled toward $\mathbf{a}^{\text{MPC}}$ at the WM-visited latents --- so the clean offline $Q$ remains the ranker and the reward clash is moot (the only consumer of rewards, the critic, never updates). This single change lifts the standalone policy to $90.7 \pm 0.8\%$ final and $94.8 \pm 0.7\%$ peak: a $+14$-point gain over the offline start, recovered with no real-environment interaction and no planner at test time.

\begin{table}[h]
\centering
\caption{Offline-to-online inside the world model, single task
   (\texttt{visual-cube-single-singletask-task1-v0}). The online phase runs $500\,000$
   FMQ-MPC steps ($N = 32$, $H = 2$, $\eta = 0.02$, Q-all) \emph{inside the WM}, adding each
   imagined transition $(z, \mathbf{a}^{\text{MPC}}, r_\psi, z'_\psi)$ to a buffer seeded with the
   offline data. Success rate is \emph{standalone} (best-of-$4$, no planner at test time), $250$
   episodes, mean $\pm$ std over $3$ seeds. E2a and E2c differ only in whether the critic is updated
   on the mixed buffer.}
\label{tab:mpc-online-e2}
\begin{tabular}{l l r r}
\toprule
Run & Critic during online & Final SR (\%) & Peak SR (\%) \\
\midrule
Offline FQL start          & --- (frozen offline)         & 76.4 & 88.0 \\
E2a                        & updated on real $+$ imagined & 30.7 $\pm$ 20.8 & 89.5 $\pm$ 0.7$^{\dagger}$ \\
E2c                        & frozen (real only)           & \textbf{90.7 $\pm$ 0.8} & \textbf{94.8 $\pm$ 0.7} \\
\midrule
E2c $+$ FMQ-MPC at test    & frozen                       & \textit{(running)} & \textit{(running)} \\
\bottomrule
\end{tabular}

\vspace{2pt}
{\footnotesize $^{\dagger}$ E2a's peak is transient --- reached near $50$k steps before the critic destabilises; by $500$k all three seeds have collapsed (final $60/14/18\%$).}
\end{table}

\paragraph{Goal-conditioned extension (E3).} The same agent extends to goal-conditioned control with no change to the RL core: the observation becomes the concatenation $[z, z_{\text{goal}}]$, the agent is trained jointly across all five \texttt{visual-cube-single} tasks, and we add hindsight relabelling \cite{her} with the HIQL-style indicator reward \cite{park2024hiqlofflinegoalconditionedrl} in place of any latent-distance reward. A pure-offline goal-conditioned checkpoint reaches a $5$-task standalone average of $90.0\%$, and the online-in-WM phase of E2c is then applied on top. Table~\ref{tab:mpc-online-e3} reports the ablation. The honest finding is that the online phase does \emph{not} improve over the strong offline baseline on final success: the best variants peak around $92\%$ but settle at $80$--$85\%$, which we attribute to coverage dilution (the buffer drifts from fully diverse offline play toward the narrow slice of latent space the MPC controller visits) compounded by the larger $(z, z_{\text{goal}})$ space the goal-conditioned policy must cover. The leakage-free ``fair'' variant, in which online rollout goals are random achieved states rather than the evaluation goals, is the strongest (E3g). One consistency check is worth highlighting: because the goal-conditioned setting uses a single indicator-reward convention everywhere, the reward clash of E2a is absent, and leaving the critic \emph{unfrozen} (E3e) no longer collapses --- direct corroboration that the E2a collapse was driven by the reward-convention mismatch and the WM-in-target mechanism rather than by the mere presence of imagined data.

\begin{table}[h]
\centering
\caption{Goal-conditioned extension (E3): multi-task across all five \texttt{visual-cube-single}
   tasks with hindsight relabelling, $[z, z_{\text{goal}}]$ observations, and the HIQL-style
   indicator reward. Online phase as in Table~\ref{tab:mpc-online-e2}; SR is the standalone $5$-task
   average ($50$ episodes/task), mean $\pm$ std over $3$ seeds unless noted. \emph{Goal source} is
   how online rollout goals are drawn: the five evaluation goals, or random achieved states (the
   leakage-free ``fair'' setting).}
\label{tab:mpc-online-e3}
\begin{tabular}{l l l l r r}
\toprule
Run & Sampling & Critic & Goal source & Final SR (\%) & Peak SR (\%) \\
\midrule
Offline GC baseline & ---        & ---     & ---            & 90.0 & 90.0 \\
\midrule
E3b & HER        & frozen  & eval goals     & 80.1 $\pm$ 1.6 & 92.1 $\pm$ 0.2 \\
E3c & no-HER     & frozen  & eval goals     & 73.2 $\pm$ 1.2 & 84.8 $\pm$ 2.0 \\
E3d & mix 50/50  & frozen  & eval goals     & 77.8 $\pm$ 3.0$^{\ddagger}$ & 87.4 $\pm$ 1.0$^{\ddagger}$ \\
E3f & mix 75/25  & frozen  & eval goals     & 80.9 $\pm$ 1.5 & 87.9 $\pm$ 0.7 \\
E3e & HER        & updated & eval goals     & 81.5 $\pm$ 1.3 & 90.9 $\pm$ 0.7 \\
\midrule
E3g       & HER       & frozen & random states & 85.0 $\pm$ 0.6$^{\S}$ & 91.0 $\pm$ 0.6$^{\S}$ \\
E3g-clean & mix 50/50 & frozen & random states & \textit{(running)} & \textit{(running)} \\
\bottomrule
\end{tabular}

\vspace{2pt}
{\footnotesize $^{\ddagger}$ Two complete seeds (the third was truncated by a disk fault at $75$k).\quad
$^{\S}$ Two seeds; the third is still running.}
\end{table}

\paragraph{Layering the planner back on (in progress).} E2c yields a strong standalone actor; the obvious last step is to apply the test-time FMQ-MPC planner (\S\ref{sec:mpc-fmq}) on top of the best E2c checkpoints, exactly as the distillation arc did to reach $96.0\%$ in Table~\ref{tab:mpc-fmq-distilled}. That evaluation is running at the time of writing; the corresponding row of Table~\ref{tab:mpc-online-e2} is left to be filled in.

\section{Discussion}
\label{sec:mpc-discussion}

\paragraph{Why this works where WM-environment training failed.} The world model is used at inference only, only over short horizons ($H \le 2$), and only as a source of \emph{rankings} between candidate actions. It is never the source of a TD target. The critic continues to be trained on real $(z, \mathbf{a}, r, z')$ from the offline buffer, so the four failure mechanisms of \S\ref{sec:oo-synthesis} --- compounding error, $Q$-corruption, threshold pathology, and model exploitation --- never come into play. The pivot is structural rather than algorithmic: we replaced the question ``can the WM substitute for the environment?'' with ``can the WM rank a small candidate set?'' and the answer to the second question is much friendlier than the first.

\paragraph{Training-inference symmetry under FMQ.} Tables~\ref{tab:mpc-fmq-basepolicy} and~\ref{tab:mpc-fmq-distilled} form a clean before/after pair. The base critic peaks under FMQ inference at $\eta \approx 0.1$ and collapses above; the FMQ-distilled critic peaks at $\eta \approx 0.02$--$0.03$ and stays above $94\%$ across $[0.01, 0.05]$. The shift is large enough that it is worth naming the mechanism explicitly: the $Q$-function calibrates to the action distribution it is supervised on, so training on FMQ-shaped actions makes the critic reliable on FMQ-shaped actions, which then permits much finer trust-region refinement at inference. This is an instance of the broader principle that the distribution of states-and-actions the critic sees at training should match the distribution the critic will be queried on at inference.

\paragraph{Reintroducing the world model, safely.} \S\ref{sec:mpc-online-wm} closes the chapter by doing the one thing the distillation methods refused to do --- putting world-model transitions back into the replay buffer --- and shows the refusal was sufficient but not necessary. What is actually necessary is narrower: the critic must never bootstrap on the imagined successors. With the offline critic frozen, it plays a dual role that keeps the whole loop on the data manifold: it \emph{selects} which imagined transitions are generated (it is the score function inside MPC) while never being \emph{corrupted} by them (it does not update). The world model is admitted only in the capacity the previous chapter exonerated it for --- ranking a small set of short rollouts --- and the standalone policy gains $+14$ points (E2c, Table~\ref{tab:mpc-online-e2}). This is the same training-inference matching principle as the FMQ symmetry above, in its strongest form: freezing the critic guarantees it is queried at inference on exactly the real-data distribution it was trained on, which is precisely why the imagined data can be trusted to be filtered correctly.

\paragraph{Limitations and next steps.} Most experiments in this chapter are on a single OGBench task; the goal-conditioned runs of \S\ref{sec:mpc-online-wm} are the exception, spanning all five \texttt{cube-single} tasks jointly, but we have not tested cross-environment generalisation, and in that multi-task setting the online-in-WM phase does not yet improve on the offline baseline. There is no real online interaction at any point; distillation is offline throughout, and the ``online'' phase of \S\ref{sec:mpc-online-wm} is online only with respect to the world model. The dominant compute cost is the periodic relabelling pass, which scales linearly in $N$, $H$, and the buffer size, and we have not explored amortising this across tasks. Finally, the success of the inference-time procedure depends on the WM being accurate enough over $1$--$2$ steps; longer planning horizons remain blocked by the compounding error of \S\ref{sec:oo-exp1}. The FMQ paper also describes a complementary inference-time refinement, $Q$-Guided Beam Search, that combines stochastic renoising with beam search over the trust region (\S\ref{sec:fmq-bg}); evaluating it on top of our distilled actors is a natural extension we leave for future work.
